% ৮.১ Compiler
\section{Compiler}

\begin{learningbox}
এই টপিক শেষে শিক্ষার্থী compiler-এর কাজ, সুবিধা, অসুবিধা এবং interpreter-এর সাথে পার্থক্য লিখতে পারবে।
\end{learningbox}

\begin{definitionbox}{Compiler}
যে translator পুরো source program একসাথে machine code/object code-এ রূপান্তর করে এবং error list দেখায়, তাকে compiler বলা হয়।
\end{definitionbox}

\begin{center}
\fbox{\parbox{0.88\textwidth}{
\centering
\textbf{Compiler flow placeholder}\\[4pt]
Source code $\rightarrow$ Compiler $\rightarrow$ Object code $\rightarrow$ Linker $\rightarrow$ Executable
}}
\end{center}

\begin{keypointbox}{সুবিধা ও সীমাবদ্ধতা}
\begin{itemize}
\item একবার compile হলে executable দ্রুত run হয়।
\item object/executable file তৈরি হতে পারে।
\item সব error একসাথে দেখাতে পারে, তাই নতুনদের কাছে debugging কঠিন লাগতে পারে।
\end{itemize}
\end{keypointbox}

\begin{examplebox}{C language}
C program সাধারণত compiler দিয়ে machine code-এ রূপান্তরিত হয়, তারপর run করা হয়।
\end{examplebox}
